%!TEX root = ../main.tex

\section{Limitations}

Despite strong retrieval performance, \retrieverShort has limitations. First, agentic retrieval incurs higher latency than single-pass retrievers because it requires multiple LLM calls over an interaction trajectory. Larger budgets improve retrieval quality but also increase runtime. Second, our best-performing configuration relies on a large proprietary LLM, which can constrain scalability due to cost and availability. While \retrieverShort is LLM-agnostic, retrieval quality can drop with smaller models; we partially mitigate this via trajectory distillation into Qwen3-8B \cite{yang_qwen3_2025}, though distilled agents still trail the teacher on challenging regimes. Third, \retrieverShort assumes that node descriptors and relation information are sufficiently informative for BM25 global search and for ranking neighborhood expansions. Sparse or templated text can prevent the agent from locating relevant seed nodes or disambiguating them. Because the global search is lexical, mismatches in vocabulary (\textit{e.g.}, paraphrases and domain-specific aliases) can cause under-retrieval. Fourth, our evaluation is centered on text-rich KG benchmarks, so performance gains may not transfer to graphs with limited text descriptions. 

Although \retrieverShort is a general retrieval approach, agentic graph exploration can create risks if used without safeguards. Retrieval errors can be treated as support for downstream decisions, and interaction traces may expose sensitive attributes if node text contains private information. Mitigation requires redaction for sensitive fields and bias audits prior to deployment.


\section{Ethical Considerations}

This study does not use human annotators, crowdworkers, or research with human participants. Ethical concerns arise in downstream use of agentic retrieval over text-rich knowledge graphs. Retrieval errors can be treated as evidence and multi-step exploration can surface sensitive attributes present in graph text.
The approach may also amplify biases in the underlying graph. If some languages and communities have sparse descriptions or different naming conventions, global lexical search and neighborhood ranking may under-retrieve relevant information, leading to unequal coverage across groups and reduced benefits for underrepresented stakeholders.
We recommend safeguards before deployment, including redaction of sensitive fields and bias audits.
Potential positive impact includes improving access to large knowledge graphs for language models, including information that may be difficult to access with text retrieval alone.